\clearpage
\chapter*{第三部分结语}
\addcontentsline{toc}{chapter}{第三部分结语}
\label{part:03-closing}

\begin{center}
  \includegraphics[width=\linewidth,height=.43\textheight,keepaspectratio]{assets/part-art/part-3-closing.png}
  \par\vspace{4pt}
  {\scriptsize\color{Muted}取意于塞尚《圣维克多山》\par}
  \BookEpigraph{part3}
\end{center}

Alice 的原购买曾经只在一台服务的事务中出现。到了这一部分，它已经随日志传到接班副本，进入快照，也可能在购买库的旧备份中消失。平台上的 X1 与履约库里的 F501 继续存在，恢复者要把它们接回原来的 P1/C1/O1。业务的过去分布在不同保存位置，每一个位置都有自己的提交与恢复边界。

复制首先延续已经确认的决定。状态机按提交顺序应用命令，把身份绑定、原输入、结果、会话资格和应用边界一起保存。另一副本接班以后，重复命令即使位于新的日志位置，也会找到原 $S$。快照压缩了日志表示，原购买及其拒绝重复的依据继续属于恢复后的状态。

记录回收改变了保存这些区别的方式。E1 关闭以后，旧购买提交入口结束准入，原来的逐项结果可以按规则缩减，关闭决定和从未重用的登记进度继续保存。已经接受的 C1/T1 则沿原关联完成核查。平台逐键窗口、旧在途请求与可取得的终结证据共同决定发送何时结束，固定的本地截止把客户端能够重新提交的范围写清楚。

状态表达也能改变需要传播的历史。标签按并集合并，重复到达后集合保持；带来源标记的删除把已经观察到的新增与并发新增区分开来。资金增量、退款额度和首次发货依赖另外的业务条件。把同一笔 1000 退款传播两次，与接受两项各 1000 的退款，在原 RF 上有不同含义。已经形成的事实可以继续合并，首次责任由身份、授权与共同仲裁决定。

\CFifteenFigure{ch15-22-part3-history}{保存决定、关闭旧准入和恢复原关系共同延续已经承担的业务责任}{fig:part03-history}

主动回退购买库时，独立 J 保存的原关系提供了另一条恢复路径。恢复者先从 $S$ 取得原 P1/C1/T1，再用原范围和 \code{pay-C1} 核对 X1，并以 $R_J$ 中的原编号、结果和 EV1 物化本地记录。已有行的完整内容参与核对，冲突使本批写入与物化边界一起回滚。目录把恢复资格与新意图准入连接起来，入口依据已经完整恢复的范围开放。

退款把这条保存链延伸到新责任。RF1 的平台结果未知时，原金额继续占用 X1 的退款额度；成功证据到达后，占用转为已退，足够最终拒绝则释放额度。旧备份删除本地投影以后，J 中的绑定与预占依然能够重建。净资金、原购买成功和当前退款进度各自有确定的数据来源，查询把这些事实组合到同一项业务上。

这些机制进入完整系统以后，会同时争用目录、数据库、任务队列和外部平台。架构设计要先给出它们共同构成的系统：谁拥有每类事实，哪些调用跨过事务域，流量如何到达权威位置，恢复者怎样发现并接续已经接受的责任。在这个整体中，再决定每个组件的数据、接口、容量与故障处理。第四部分由此转向完整系统架构，让购买、付款、退款和履约的保证在同一部署中各有承担位置。
\FloatBarrier
